\newcommand{\ReshapingNodeFigure}{%
\begin{figure}[H]
\centering
\begin{tikzpicture}[x=.88cm,y=2.8cm,font=\small,
  polyshape/.style={draw=teal!65!black,fill=teal!10,line width=.8pt},
  slice/.style={draw=orange!85!black,line width=1.2pt}]
  \begin{scope}
    \node at (1.4,.77) {$Q$};
    \draw[->,gray!65] (-.2,0)--(3.35,0);
    \draw[->,gray!65] (0,-.04)--(0,.63);
    \path[polyshape] (0,0)--(3,0)--(.5,.5)--(0,.375)--cycle;
    \node[below left] at (0,0) {$0$};
    \node[below] at (3,0) {$(a,0)$};
    \node[above right] at (.5,.5) {$(c,c)$};
    \node[left] at (0,.375) {$(0,b)$};
    \draw[slice] (0,.18)--(.26,.44);
    \node at (1.4,-.21) {diagonal slices};
  \end{scope}
  \begin{scope}[xshift=4.65cm]
    \node at (1.35,.77) {after the shear};
    \draw[->,gray!65] (-.7,0)--(3.35,0);
    \draw[->,gray!65] (0,-.04)--(0,.63);
    \path[polyshape] (0,0)--(3,0)--(0,.5)--(-.375,.375)--cycle;
    \node[below right] at (0,0) {$0$};
    \node[below] at (3,0) {$(a,0)$};
    \node[above right] at (0,.5) {$(0,c)$};
    \node[left] at (-.375,.375) {$(-b,b)$};
    \draw[slice] (-.18,.18)--(-.18,.44);
    \node at (1.35,-.21) {vertical slices};
  \end{scope}
  \begin{scope}[xshift=9.3cm]
    \node at (1.35,.77) {$P$};
    \draw[->,gray!65] (-.7,0)--(3.35,0);
    \draw[->,gray!65] (0,-.04)--(0,.63);
    \path[polyshape] (-.375,0)--(3,0)--(0,.5)--cycle;
    \node[below left] at (-.375,0) {$(-b,0)$};
    \node[below] at (3,0) {$(a,0)$};
    \node[above right] at (0,.5) {$(0,c)$};
    \draw[slice] (-.18,0)--(-.18,.26);
    \node at (1.35,-.21) {slices moved to height $0$};
  \end{scope}
  \draw[-{Stealth[length=2mm]},thick] (3.65,.23)--(4.25,.23);
  \draw[-{Stealth[length=2mm]},thick] (8.9,.23)--(9.5,.23);
\end{tikzpicture}
\caption{Bounding polygons in the diagonal compression.  The shear turns
lines of direction $(1,1)$ into vertical lines.  Moving each vertical slice
to start at height zero produces $P$; the marked slices correspond under
these operations.  These are continuous bounds: actual row exponents are
replaced by possible vanishing orders, not translated point by point.
All three polygons have area $H/2$.  Axis scales are chosen for legibility.}
\label{fig:node-compression}
\end{figure}%
}

\newcommand{\ReshapingCompressionFigure}{%
\begin{figure}[H]
\centering
\begin{tikzpicture}[x=.88cm,y=.88cm,font=\small,
  polyshape/.style={draw=teal!65!black,fill=teal!10,line width=.8pt},
  slice/.style={draw=orange!85!black,line width=1.2pt}]
  \begin{scope}
    \node at (1.5,3.25) {in the basis $(u,u')$};
    \path[polyshape] (1,0)--(3,.9)--(.4,2.4)--cycle;
    \draw[densely dashed,gray] (-.1,0)--(1,0);
    \draw[densely dashed,gray] (-.1,.9)--(.775,.9);
    \draw[densely dashed,gray] (-.1,2.4)--(.4,2.4);
    \node[left] at (-.1,0) {$q_0$};
    \node[left] at (-.1,.9) {$q_1$};
    \node[left] at (-.1,2.4) {$q_2$};
    \draw[slice] (.775,.9)--(3,.9);
    \node[above] at (1.9,.9) {$w$};
    \node at (1.5,-.62) {height range $R$};
  \end{scope}
  \begin{scope}[xshift=4.65cm]
    \node at (1.35,3.25) {horizontal compression};
    \draw[->,gray!65] (0,-.12)--(0,2.75);
    \path[polyshape] (0,0)--(2.225,.9)--(0,2.4)--cycle;
    \node[left] at (0,0) {$q_0$};
    \node[left] at (0,.9) {$q_1$};
    \node[left] at (0,2.4) {$q_2$};
    \draw[slice] (0,.9)--(2.225,.9);
    \node[right] at (2.225,.9) {$(w,q_1)$};
    \draw[gray!70,densely dashed] (.85,.344)--(.85,1.827);
    \node at (1.35,-.62) {maximum slice length $w$};
  \end{scope}
  \begin{scope}[xshift=9.3cm]
    \node at (1.4,3.25) {vertical compression};
    \draw[->,gray!65] (-.12,0)--(2.65,0);
    \draw[->,gray!65] (0,-.12)--(0,2.75);
    \path[polyshape] (0,0)--(2.225,0)--(0,2.4)--cycle;
    \node[below left] at (0,0) {$0$};
    \node[below] at (2.225,0) {$w$};
    \node[left] at (0,2.4) {$R$};
    \draw[slice] (.85,0)--(.85,1.483);
    \node at (1.4,-.62) {$\Delta(w,R)$};
  \end{scope}
  \draw[-{Stealth[length=2mm]},thick] (3.55,1.65)--(4.15,1.65);
  \draw[-{Stealth[length=2mm]},thick] (8.75,1.65)--(9.35,1.65);
\end{tikzpicture}
\caption{The primitive direction is chosen so that the transverse height
range is $R=H/w$.  Since the area is $H/2$, the longest horizontal slice
has length $w$.  Horizontal compression aligns the left endpoints of
these slices.  Vertical compression then gives the right triangle with
intercepts $w$ and $R$.  The middle and right panels mark vertical slices
of equal length; the drawing is schematic.}
\label{fig:final-compressions}
\end{figure}%
}
